\scanpage{19}
On obtient ainsi, en termes de \deleted{\unsure{$O$, un ``g.}\ \unsure{pur}} \inserted{$C$, une} Gr-catégorie, qui, pour un objet $X$ fixé de $C$ d'anneau $O$, (i.e. une réalisation de $C$ comme $\unsure{\mathbb P|P}(O)$) s'écrit comme la \deleted{\unsure{groupoïde} de Brand} \inserted{catégorie des} bimodules ``bi-inversibles''. Il lui correspond les invariants
\[
\underset{\substack{|\\\text{groupe de Brand}}}{\operatorname{Brd}(O)}
=\pi_{\deleted{\unsure{0}}1}(\uline{\operatorname{Aut}}(C))
\]
\[
O^{\unsure{\natural|\sharp}*}=\pi_0(\uline{\operatorname{End}}(C))=\operatorname{End}(\operatorname{id}_C)
\]
et \deleted{une classe canonique} l'opération \inserted{\uline{triviale}} de $\operatorname{Brd}(O)$ sur $O^{\unsure{\natural|\sharp}*}$ \deleted{\unsure{et}} \deleted{$\xi\in H^3(\operatorname{Brd}$} [l'image \inserted{$\sigma(\lambda)$} par $\unsure{\lambda|l}=\operatorname{cl}(L)\in\operatorname{Brd}(O)$ \deleted{de l'élé} ($L$ bimodule biinversible) \deleted{\unsure{étant définie}} \inserted{\unsure{conditionnée}} par la condition
\[
\sigma(\lambda)x=x\lambda\qquad\text{pour tt }x\in L,\ \lambda\in O^{\unsure{\natural|\sharp}}
\]
i.e., \unsure{moralement}, $\sigma(\lambda)=x\lambda x^{-1}=\lambda$ car $\lambda$ central --]

et un élément canonique
\[
\xi\in H^3(\operatorname{Brd}(O),O^{\unsure{\natural|\sharp}*})
\]

Les catégories $C$ de Deuring équivalentes entre elles forment une \deleted{2-catégorie} 2-groupoïde \inserted{$\uuline{\mathcal D}$}, dont les \deleted{isoms} Gr-catégories fondamentales sont les groupoïdes de Brandt des $C$ envisagées, i.e. des $O$ envisagés. La deux-catégorie est aussi 2-équivalente à la catégorie \deleted{\unsure{\mathcal D}} \inserted{$\uline{\operatorname{Brd}}$} dont les \uline{objets} sont les $O$ \deleted{\unsure{apparentés}\ \unsure{à une classe de Deuring d'anneaux}} \inserted{d'une classe de Deuring d'} anneaux, \deleted{\unsure{les hom}} la catégorie des hom de $O$ dans $O'$ est la catégorie des bitorseurs biinversibles sous $O'$, \unsure{et|$O$}, la composition des hom est le $\otimes$ \ldots
\mspage{18}
